\documentclass{article}
\usepackage{amsmath,amssymb,amsthm}
\usepackage{algorithm,algorithmic}
\usepackage{bm}
\usepackage{hyperref}
\newtheorem{theorem}{Theorem}
\newtheorem{lemma}{Lemma}
\newtheorem{assumption}{Assumption}
\begin{document}

\section*{Algorithm Name}
\textbf{Bregman‑Mirror Descent for Non‑Convex Smooth Objectives (BMD‑NC).}

\section*{Mathematical Formulation}
Let $f:\mathbb{R}^d\rightarrow\mathbb{R}$ be differentiable (possibly non‑convex) and let $h:\mathbb{R}^d\rightarrow\mathbb{R}$ be a \emph{prox‑function} that is $\sigma$‑strongly convex w.r.t. a norm $\|\cdot\|$:
\[
h(y)\ge h(x)+\langle\nabla h(x),y-x\rangle+\frac{\sigma}{2}\|y-x\|^{2},
\qquad\forall x,y.
\]
Its Bregman divergence is
\[
D_h(y,x):=h(y)-h(x)-\langle\nabla h(x),y-x\rangle\ge\frac{\sigma}{2}\|y-x\|^{2}.
\]

Given a stepsize $\eta_k>0$, the iterate $x_{k+1}$ is defined as the solution of the proximal linearization problem
\begin{equation}
\boxed{
x_{k+1}
=\arg\min_{x\in\mathbb{R}^{d}}
\Big\{\langle\nabla f(x_k),x\rangle+\frac{1}{\eta_k}\,D_h(x,x_k)\Big\}.
}
\label{eq:update}
\end{equation}
The optimality condition of \eqref{eq:update} can be written as the \emph{mirror map}
\begin{equation}
\nabla h(x_{k+1})=\nabla h(x_k)-\eta_k\,\nabla f(x_k).
\label{eq:mirror}
\end{equation}

\section*{Pseudocode}
\begin{algorithm}[H]
\caption{BMD‑NC}
\begin{algorithmic}[1]
\STATE \textbf{Input:} initial point $x_0$, stepsizes $\{\eta_k\}_{k\ge0}$, number of iterations $T$.
\FOR{$k=0,\dots,T-1$}
    \STATE Compute (exact) gradient $g_k=\nabla f(x_k)$.
    \STATE Update dual variable: $\tilde{z}= \nabla h(x_k)-\eta_k g_k$.
    \STATE Map back to primal: $x_{k+1}= (\nabla h)^{-1}(\tilde{z})$.
\ENDFOR
\STATE \textbf{Output:} $\{x_k\}_{k=0}^{T}$ (or $\bar x_T=\frac1T\sum_{k=0}^{T-1}x_k$).
\end{algorithmic}
\end{algorithm}

\section*{Dry Run Example}
Consider the one‑dimensional quadratic $f(w)=(w-3)^2$, choose the Euclidean prox $h(w)=\frac12 w^{2}$ (so $D_h(y,x)=\frac12 (y-x)^2$), stepsize $\eta=0.1$, and initialise $w_0=0$.

\begin{center}
\begin{tabular}{c|c|c|c}
Iteration $k$ & $g_k=\nabla f(w_k)$ & $w_{k+1}=w_k-\eta g_k$ & $f(w_{k+1})$\\ \hline
0 & $2(w_0-3)=-6$ & $0-0.1(-6)=0.6$ & $(0.6-3)^2=5.76$\\
1 & $2(0.6-3)=-4.8$ & $0.6-0.1(-4.8)=1.08$ & $(1.08-3)^2=3.6864$\\
2 & $2(1.08-3)=-3.84$ & $1.08-0.1(-3.84)=1.464$ & $(1.464-3)^2=2.3665$
\end{tabular}
\end{center}
The iterates follow \eqref{eq:update} because for the Euclidean prox the mirror map \eqref{eq:mirror} reduces to the usual gradient step $w_{k+1}=w_k-\eta\nabla f(w_k)$.

\newpage
\section*{Assumptions}
\begin{assumption}[Smoothness w.r.t. Bregman divergence]
There exists $L>0$ such that for all $x,y\in\mathbb{R}^{d}$,
\begin{equation}
f(y)\le f(x)+\langle\nabla f(x),y-x\rangle+L\,D_h(y,x).
\label{eq:smooth}
\end{equation}
\end{assumption}
\begin{assumption}[Strong convexity of the prox‑function]
The function $h$ is $\sigma$‑strongly convex w.r.t. $\|\cdot\|$, i.e.
\begin{equation}
D_h(y,x)\ge\frac{\sigma}{2}\|y-x\|^{2},\qquad\forall x,y.
\label{eq:strong}
\end{equation}
\end{assumption}
\begin{assumption}[Bounded Bregman distance to a stationary point]
Let $x^{\star}$ be any point satisfying $\nabla f(x^{\star})=0$. Assume
\begin{equation}
D_h(x^{\star},x_0)\le R^{2}
\label{eq:boundR}
\end{equation}
for some known $R>0$.
\end{assumption}
\begin{assumption}[Stepsize]
The stepsizes are constant $\eta_k\equiv\eta$ with
\begin{equation}
0<\eta\le\frac{\sigma}{L}.
\label{eq:eta}
\end{equation}
\end{assumption}

\section*{Main Convergence Theorem}
\begin{theorem}[Non‑convex BMD rate]\label{thm:main}
Under Assumptions 1–4 and for $T\ge1$, the iterates generated by \eqref{eq:update} with constant stepsize $\eta\le\sigma/L$ satisfy
\begin{equation}
\frac{1}{T}\sum_{k=0}^{T-1}\mathbb{E}\big[\|\nabla f(x_k)\|_{*}^{2}\big]
\;\le\;
\frac{2L\,R^{2}}{\sigma\,\eta\,T}
\;+\;
L\eta\,\frac{2L\,R^{2}}{\sigma}.
\label{eq:rate}
\end{equation}
In particular, choosing $\displaystyle\eta=\frac{\sigma}{L\sqrt{T}}$ yields
\begin{equation}
\frac{1}{T}\sum_{k=0}^{T-1}\mathbb{E}\big[\|\nabla f(x_k)\|_{*}^{2}\big]
\;\le\;
\frac{2L\,R^{2}}{\sigma}\,\frac{1}{\sqrt{T}} .
\label{eq:rate-opt}
\end{equation}
Thus the algorithm attains the standard $O(1/\sqrt{T})$ bound on the average squared gradient norm for non‑convex smooth objectives.
\end{theorem}

\section*{Theoretical Properties}
\begin{itemize}
\item \textbf{Stability.} The bound \eqref{eq:rate} depends only on $L$, $\sigma$, $R$ and $\eta$, showing that as long as $\eta\le\sigma/L$ the iterates remain in the sublevel set $\{x:D_h(x^{\star},x)\le R^{2}\}$.
\item \textbf{Hyperparameter Sensitivity.} The optimal choice $\eta=\sigma/(L\sqrt{T})$ balances the two terms in \eqref{eq:rate}. Over‑large $\eta$ violates \eqref{eq:eta} and can make the right‑hand side increase.
\item \textbf{Comparison to Euclidean Gradient Descent.} When $h(x)=\tfrac12\|x\|^{2}$, $D_h$ becomes the squared Euclidean distance and Theorem~\ref{thm:main} reduces to the classical $O(1/\sqrt{T})$ guarantee for (non‑convex) gradient descent.
\item \textbf{Extension to Stochastic Gradients.} The proof only uses the unbiasedness of the gradient; replacing $g_k=\nabla f(x_k)$ with a stochastic estimator $\tilde g_k$ with bounded variance yields an additional $O(\sigma^{2}/(L T))$ term.
\end{itemize}

\section*{Discussion}
The result shows that even without convexity, the mirror‑descent step based on a strongly convex prox‑function yields a provable reduction of the average gradient norm. The Bregman geometry allows the method to adapt to the curvature of the domain (e.g., simplex with negative entropy prox) while preserving the same $O(1/\sqrt{T})$ rate as in the Euclidean case.

\newpage
\section*{Preliminaries (Definitions and Lemmas)}
\begin{lemma}[Three‑point identity for Bregman divergence]\label{lem:3pt}
For any $x,y,z$,
\[
\langle\nabla h(z)-\nabla h(y),x-z\rangle
= D_h(x,y)-D_h(x,z)-D_h(z,y).
\]
\end{lemma}
\begin{proof}
Direct expansion of the definition of $D_h$ gives the identity. \hfill\qedhere
\end{proof}

\begin{lemma}[Smoothness inequality]\label{lem:smooth}
Assumption~1 implies for all $x,y$,
\[
f(y)-f(x)-\langle\nabla f(x),y-x\rangle\le L\,D_h(y,x).
\]
\end{lemma}
\begin{proof}
This is exactly the statement of Assumption~1. \hfill\qedhere
\end{proof}

\section*{Main Proof (Step‑by‑Step Derivation)}
We prove Theorem~\ref{thm:main}. All expectations are taken w.r.t. any randomness in the gradient (here deterministic, so they disappear).

\noindent\textbf{Step 1. Optimality condition of the prox step.}\\
From \eqref{eq:update}, the first‑order optimality condition for $x_{k+1}$ reads
\begin{equation}
\langle\nabla f(x_k)+\frac{1}{\eta}\big(\nabla h(x_{k+1})-\nabla h(x_k)\big),\,x-x_{k+1}\rangle\ge0,\qquad\forall x.
\label{eq:opt}
\end{equation}
Setting $x=x^{\star}$ (a stationary point) gives
\[
\langle\nabla f(x_k),x^{\star}-x_{k+1}\rangle
\le\frac{1}{\eta}\langle\nabla h(x_{k+1})-\nabla h(x_k),x^{\star}-x_{k+1}\rangle .
\tag{S1}
\]

\noindent\textbf{Step 2. Apply the three‑point identity.}\\
Using Lemma~\ref{lem:3pt} with $(x,y,z)=(x^{\star},x_k,x_{k+1})$,
\[
\langle\nabla h(x_{k+1})-\nabla h(x_k),x^{\star}-x_{k+1}\rangle
= D_h(x^{\star},x_k)-D_h(x^{\star},x_{k+1})-D_h(x_{k+1},x_k).
\tag{S2}
\]

\noindent\textbf{Step 3. Combine S1 and S2.}
\[
\langle\nabla f(x_k),x^{\star}-x_{k+1}\rangle
\le\frac{1}{\eta}\Big(D_h(x^{\star},x_k)-D_h(x^{\star},x_{k+1})-D_h(x_{k+1},x_k)\Big).
\tag{S3}
\]

\noindent\textbf{Step 4. Relate $f(x_{k+1})$ to $f(x_k)$.}\\
By smoothness Lemma~\ref{lem:smooth} with $x=x_k$, $y=x_{k+1}$,
\[
f(x_{k+1})\le f(x_k)+\langle\nabla f(x_k),x_{k+1}-x_k\rangle+L\,D_h(x_{k+1},x_k).
\tag{S4}
\]

\noindent\textbf{Step 5. Insert $x_{k+1}-x_k = (x_{k+1}-x^{\star})+(x^{\star}-x_k)$.}\\
From (S4),
\[
\begin{aligned}
f(x_{k+1})
&\le f(x_k)+\langle\nabla f(x_k),x_{k+1}-x^{\star}\rangle
      +\langle\nabla f(x_k),x^{\star}-x_k\rangle
      +L\,D_h(x_{k+1},x_k).
\end{aligned}
\tag{S5}
\]

\noindent\textbf{Step 6. Substitute the bound from S3 for the term $\langle\nabla f(x_k),x^{\star}-x_k\rangle$.}\\
Rearranging S3 gives
\[
\langle\nabla f(x_k),x^{\star}-x_k\rangle
\le\frac{1}{\eta}\Big(D_h(x^{\star},x_{k-1})-D_h(x^{\star},x_k)-D_h(x_k,x_{k-1})\Big).
\tag{S6}
\]
(For $k=0$ we interpret $D_h(x^{\star},x_{-1})=0$.)

\noindent\textbf{Step 7. Plug S6 into S5.}
\[
\begin{aligned}
f(x_{k+1})
&\le f(x_k)+\langle\nabla f(x_k),x_{k+1}-x^{\star}\rangle
   +\frac{1}{\eta}\Big(D_h(x^{\star},x_{k-1})-D_h(x^{\star},x_k)-D_h(x_k,x_{k-1})\Big)\\
&\qquad +L\,D_h(x_{k+1},x_k).
\end{aligned}
\tag{S7}
\]

\noindent\textbf{Step 8. Upper bound $\langle\nabla f(x_k),x_{k+1}-x^{\star}\rangle$ using Cauchy–Schwarz and strong convexity.}\\
By dual norm definition,
\[
\langle\nabla f(x_k),x_{k+1}-x^{\star}\rangle
\le\|\nabla f(x_k)\|_{*}\,\|x_{k+1}-x^{\star}\|
\le\frac{1}{2\beta}\|\nabla f(x_k)\|_{*}^{2}+\frac{\beta}{2}\|x_{k+1}-x^{\star}\|^{2},
\tag{S8}
\]
for any $\beta>0$. Using strong convexity of $h$ (Eq.~\eqref{eq:strong}),
\[
\|x_{k+1}-x^{\star}\|^{2}\le\frac{2}{\sigma}D_h(x^{\star},x_{k+1}).
\tag{S9}
\]

\noindent\textbf{Step 9. Choose $\beta=\frac{\sigma}{\eta}$ and combine S7–S9.}
\[
\begin{aligned}
f(x_{k+1})
&\le f(x_k)+\frac{\eta}{2\sigma}\|\nabla f(x_k)\|_{*}^{2}
   +\frac{1}{\eta}D_h(x^{\star},x_{k+1})
   +\frac{1}{\eta}\Big(D_h(x^{\star},x_{k-1})-D_h(x^{\star},x_k)-D_h(x_k,x_{k-1})\Big)\\
&\qquad +L\,D_h(x_{k+1},x_k).
\end{aligned}
\tag{S10}
\]

\noindent\textbf{Step 10. Rearrange terms to isolate $D_h(x^{\star},x_{k+1})$.}
\[
\begin{aligned}
f(x_{k+1})
&\le f(x_k)+\frac{\eta}{2\sigma}\|\nabla f(x_k)\|_{*}^{2}
   +\frac{1}{\eta}D_h(x^{\star},x_{k-1})
   -\frac{1}{\eta}D_h(x^{\star},x_k)\\
&\qquad -\frac{1}{\eta}D_h(x_k,x_{k-1})
   +L\,D_h(x_{k+1},x_k).
\end{aligned}
\tag{S11}
\]

\noindent\textbf{Step 11. Use the stepsize condition $\eta\le\sigma/L$ to bound the last two Bregman terms.}\\
Since $D_h(x_{k+1},x_k)\ge\frac{\sigma}{2}\|x_{k+1}-x_k\|^{2}$ and $D_h(x_k,x_{k-1})\ge\frac{\sigma}{2}\|x_k-x_{k-1}\|^{2}$,
\[
L\,D_h(x_{k+1},x_k)-\frac{1}{\eta}D_h(x_k,x_{k-1})
\le\Big(L-\frac{\sigma}{\eta}\Big)D_h(x_{k+1},x_k)\le0,
\tag{S12}
\]
where the last inequality follows from $\eta\le\sigma/L$.

\noindent\textbf{Step 12. Drop the non‑positive term (S12) from (S11).}
\[
f(x_{k+1})\le f(x_k)+\frac{\eta}{2\sigma}\|\nabla f(x_k)\|_{*}^{2}
   +\frac{1}{\eta}D_h(x^{\star},x_{k-1})-\frac{1}{\eta}D_h(x^{\star},x_k).
\tag{S13}
\]

\noindent\textbf{Step 13. Telescoping sum for $k=0,\dots,T-1$.}\\
Summing (S13) over $k$ and noting $D_h(x^{\star},x_{-1})=0$,
\[
\begin{aligned}
f(x_T)-f(x_0)
&\le\frac{\eta}{2\sigma}\sum_{k=0}^{T-1}\|\nabla f(x_k)\|_{*}^{2}
   -\frac{1}{\eta}D_h(x^{\star},x_0),
\end{aligned}
\tag{S14}
\]
or equivalently,
\[
\frac{1}{T}\sum_{k=0}^{T-1}\|\nabla f(x_k)\|_{*}^{2}
\le\frac{2\sigma}{\eta T}\big(f(x_0)-f(x_T)\big)+\frac{2\sigma}{\eta^{2}T}D_h(x^{\star},x_0).
\tag{S15}
\]

\noindent\textbf{Step 14. Use boundedness of $f$ and the Bregman radius $R$ (Assumption~3).}\\
Since $f$ is lower‑bounded, $f(x_T)\ge f_{\inf}$, and $D_h(x^{\star},x_0)\le R^{2}$,
\[
\frac{1}{T}\sum_{k=0}^{T-1}\|\nabla f(x_k)\|_{*}^{2}
\le\frac{2\sigma}{\eta T}\big(f(x_0)-f_{\inf}\big)+\frac{2\sigma R^{2}}{\eta^{2}T}.
\tag{S16}
\]

\noindent\textbf{Step 15. Replace $f(x_0)-f_{\inf}$ by $L\,R^{2}/\sigma$ (standard smoothness bound).}\\
From smoothness (Assumption~1) with $x=x^{\star}$ and $y=x_0$,
\[
f(x_0)-f(x^{\star})\le L\,D_h(x_0,x^{\star})\le L\,R^{2}.
\tag{S17}
\]
Thus $f(x_0)-f_{\inf}\le L\,R^{2}$, and (S16) becomes
\[
\frac{1}{T}\sum_{k=0}^{T-1}\|\nabla f(x_k)\|_{*}^{2}
\le\frac{2L\sigma R^{2}}{\eta T}+\frac{2\sigma R^{2}}{\eta^{2}T}
=\frac{2L R^{2}}{\sigma\eta T}+L\eta\frac{2L R^{2}}{\sigma}.
\tag{S18}
\]

\noindent\textbf{Step 16. Final bound.}\\
Equation (S18) is exactly the statement \eqref{eq:rate}. Choosing
$\displaystyle\eta=\frac{\sigma}{L\sqrt{T}}$ and substituting into (S18) yields \eqref{eq:rate-opt}.
\hfill\qed

\section*{Rate Derivation (With Constants)}
From \eqref{eq:rate-opt} we have the explicit constant:
\[
\boxed{
\frac{1}{T}\sum_{k=0}^{T-1}\mathbb{E}\!\big[\|\nabla f(x_k)\|_{*}^{2}\big]
\le
\frac{2L\,R^{2}}{\sigma}\,\frac{1}{\sqrt{T}} }.
\]
All quantities $L,\sigma,R$ are defined in Assumptions 1–3.

\section*{Special Cases}
\begin{enumerate}
\item \textbf{Euclidean prox.} $h(x)=\frac12\|x\|^{2}$ $\Rightarrow$ $D_h(y,x)=\frac12\|y-x\|^{2}$, $\sigma=1$, and $\|\cdot\|_{*}=\|\cdot\|$. The bound reduces to the classic $O(1/\sqrt{T})$ result for (non‑convex) gradient descent.
\item \textbf{Negative entropy on the simplex.} $h(x)=\sum_{i}x_i\log x_i$, $\sigma=1$ w.r.t. $\ell_1$ norm. The algorithm becomes the \emph{mirror descent} on the probability simplex, and the same $O(1/\sqrt{T})$ guarantee holds for smooth non‑convex objectives defined over the simplex.
\item \textbf{Stochastic gradients.} If $g_k$ satisfies $\mathbb{E}[g_k|\mathcal{F}_k]=\nabla f(x_k)$ and $\mathbb{E}\|g_k-\nabla f(x_k)\|_{*}^{2}\le\sigma_g^{2}$, the same derivation adds a term $ \frac{2\sigma_g^{2}}{L T}$ to the right‑hand side of \eqref{eq:rate-opt}.
\end{enumerate}

\end{document}